% ============================================================
% Capitolo 5 — Risultati (traduzione italiana di
% chapters/ch05_results.tex; label invariati; include le correzioni
% post GNN-B a §5.3: quattro celle, 137–142K / 162–175K)
% Ogni numero ricalcolato da episodes.jsonl su disco
% (docs/thesis/data/campaign_consolidation.json).
% ============================================================
\chapter{Risultati}
\label{ch:results}

Questo capitolo riporta i risultati della campagna in nove viste: il
pilota di aprile (\cref{sec:pilot}), integrità e setup
(\cref{sec:integrity}), efficienza campionaria
(\cref{sec:sample-efficiency}), prestazione finale
(\cref{sec:final-performance}), comportamento entro la run
(\cref{sec:behavior}), generalizzazione di mappa
(\cref{sec:generalization}), l'asse architetturale
(\cref{sec:arch-axis}), i limiti strutturali dei pedoni
(\cref{sec:ped-limits}) e i verdetti rispetto alle ipotesi
(\cref{sec:findings}).

% Fonti: ricalcolato da disco 2026-07-14 (9 run × episodes.jsonl,
% dedup + solo route_complete; integrità 6/6, 0 righe malformate).
% Figura: docs/thesis/scripts/make_pilot_collapse_fig.py.
\section{Pilota architetturale preliminare (aprile 2026)}
\label{sec:pilot}

Prima delle correzioni di correttezza della \cref{sec:bugs}, nove run
(aprile--inizio maggio 2026) hanno incrociato encoder del critic e
regimi di addestramento nell'ambiente di prima generazione: seed 42,
budget di $1$M o $3$M passi, osservazioni veicolo dell'era 44D e il
route planner pre-correzione. Sono riportate qui come studio
preliminare --- non perché i loro confronti reggano (non reggono), ma
perché hanno motivato la metodologia del \cref{ch:methodology} e
forniscono l'esempio cautelativo più chiaro di questa tesi sulle
metriche aggregate. Tutti e nove i log superano le verifiche di
integrità della \cref{sec:metrics}; la \cref{tab:pilot} riporta le
metriche cumulative ricalcolate.

\begin{table}[t]
  \centering\small
  \caption{Run del pilota di aprile 2026, ricalcolate da disco
  (cumulative, a livello di agente; \%). Tutte le run: seed 42, Town03,
  ambiente pre-correzione della lunghezza dei percorsi. La velocità per
  record non era ancora loggata in quest'era. $^{a}$~architettura non
  etichettata nei metadati della run (vecchio schema di configurazione);
  confermata MLP dall'autore. $^{b}$~variante con encoder del critic ad
  attenzione e PopArt abilitati. $^{c}$~fermata a $\approx$2.7M di 3M.}
  \label{tab:pilot}
  \begin{tabular}{@{}llllrrrrrr@{}}
    \toprule
    Run & Regime & Critic & Passi & Ep. & \multicolumn{4}{c}{Veicoli} & Ped. \\
    \cmidrule(lr){6-9}\cmidrule(l){10-10}
    & & & & & SR & s$+$t & coll. & off. & SR \\
    \midrule
    0406 & curr.\ & MLP & 1M & 1062 & 1.2 & 46.4 & 31.8 & 20.7 & 82.4 \\
    0411 & curr.\ & MLP & 3M & 3047 & 2.0 & 53.2 & 30.6 & 14.2 & 67.4 \\
    0425 & curr.\ & GNN & 3M & 3015 & 1.2 & 70.8 & 19.5 & 8.6 & 70.6 \\
    0430 & curr.\ & GNN$^{b}$ & 3M & 2188 & 1.8 & 50.9 & 28.9 & 18.5 & 66.2 \\
    0504 & curr.\ & GNN & 3M$^{c}$ & 2730 & 3.9 & 64.3 & 21.1 & 10.8 & 75.4 \\
    0405 & batch & MLP & 1M & 1066 & 1.9 & 42.3 & 38.3 & 17.6 & 77.0 \\
    0408 & batch & MLP$^{a}$ & 3M & 3038 & 0.2 & 57.8 & 28.2 & 13.8 & 72.5 \\
    0421 & batch & GNN & 3M & 3078 & 2.7 & 49.5 & 30.5 & 17.4 & 74.4 \\
    0502 & batch & GNN & 3M & 3054 & 3.0 & 58.2 & 23.6 & 15.3 & 74.4 \\
    \bottomrule
  \end{tabular}
\end{table}

Quattro osservazioni, in ordine decrescente di certezza:

\begin{enumerate}
  \item \textbf{I veicoli sedevano al pavimento in ogni cella} (SR
        0.2--3.9\%). Al pavimento, la metrica di successo non ha spazio
        per esprimere differenze: qualunque effetto architettura o
        regime possano aver avuto è compresso sotto il rumore. Il
        confronto architetturale del pilota sui veicoli è quindi
        \emph{non informativo} --- il controesempio che motiva
        l'esecuzione della campagna principale nel regime di fascia
        media. Il pavimento stesso è stato poi ricondotto al difetto di
        lunghezza dei percorsi della \cref{sec:bugs}.
  \item \textbf{I pedoni hanno appreso} (SR 66--82\%): la piattaforma
        era addestrabile per la classe più semplice anche in
        quest'era, localizzando il fallimento nel task veicolo anziché
        nello stack --- e illustrando perché tutto il reporting è per
        classe (\cref{sec:metrics}).
  \item \textbf{Una firma architetturale qualitativa.} Tra le run
        curriculum, la run con critic GNN semplice (0425) mostra il
        profilo di fallimento più difensivo delle nove: il più alto
        stuck$+$timeout (70.8\%) con la collisione più bassa (19.5\%) e
        l'offroad più basso (8.6\%), contro 53.2/30.6/14.2\% della
        controparte MLP (0411). A SR di pavimento questo è al più
        generatore di ipotesi; se una firma architetturale dipendente
        dal regime ricompaia sul trunk sano è esattamente la domanda
        H3, ripresa nella \cref{sec:arch-axis}.
  \item \textbf{Gli aggregati cumulativi possono nascondere un
        collasso.} La run 0504 è stata segnalata durante
        l'addestramento come collassata a metà corsa --- eppure il suo
        SR veicoli \emph{cumulativo} (3.9\%) è il più alto di tutte e
        nove. La \cref{fig:pilot-collapse} risolve il paradosso: la run
        ha sostenuto un SR a finestra del 10--14\% per circa 500
        episodi, poi è crollata bruscamente quasi a zero attorno
        all'episodio 1{.}100 senza mai riprendersi; il dato cumulativo è
        interamente ereditato dalla fase pre-collasso. Questa lezione
        --- la competenza corrente va misurata su una finestra, non
        dalla nascita --- ha plasmato direttamente la metrica di sblocco
        a finestra della \cref{sec:regimes} e il reporting per quartili
        della \cref{sec:metrics}.
\end{enumerate}

\begin{figure}[t]
  \centering
  \includegraphics[width=0.85\textwidth]{pilot_collapse_timeseries}
  \caption{SR veicoli a finestra (finestra di 100 episodi) per le due
  run pilota GNN-curriculum semplici. La run 0504 (tratteggiata)
  sostiene il 10--14\% per $\approx$500 episodi, collassa attorno
  all'episodio 1{.}100 e non si riprende mai --- pur chiudendo con l'SR
  cumulativo \emph{più alto} del pilota (3.9\%, \cref{tab:pilot}). Gli
  aggregati cumulativi ereditano la storia; il rilevamento dei collassi
  richiede viste a finestra.}
  \label{fig:pilot-collapse}
\end{figure}

% Fonti per §5.2–5.4, §5.6: docs/thesis/data/campaign_consolidation.json,
% ricalcolato da disco 2026-07-15 da docs/thesis/scripts/
% consolidate_campaign.py (dedup, solo route_complete, per classe).
\section{Integrità della campagna e riepilogo del setup}
\label{sec:integrity}

La \cref{tab:runmatrix} elenca le run della campagna. Tutte sono state
lanciate sullo stesso trunk (commit \texttt{7defb03}) con seed 999; i
diff di configurazione confermano che le celle differiscono solo nei
flag di regime ed encoder del critic (\cref{sec:campaign}). Tutti e
quattro i log di addestramento e tutti e quattro i log di valutazione
superano ogni condizione di integrità della \cref{sec:metrics}:
esattamente sei record per episodio ovunque, zero righe malformate, zero
valori non finiti. Un'anomalia benigna è a registro: la valutazione
GNN-curriculum ha mantenuto 401 episodi dopo la deduplica (un episodio
easy in più), uno spostamento di peso dello $0.25$\% senza effetto alla
precisione riportata.

\begin{table}[t]
  \centering\small
  \caption{Matrice delle run di campagna (tutte: 3M passi, Town03, seed
  999, trunk \texttt{7defb03}). L'esposizione di livello realizzata è
  riportata come quote easy/medium/hard degli episodi.}
  \label{tab:runmatrix}
  \begin{tabular}{@{}llrrlll@{}}
    \toprule
    Cella & Run ID & Episodi & Record & Mix di livelli (\%) & Integrità & Eval finale \\
    \midrule
    MLP-curr.\ & 20260622\_171626 & 3220 & 19{.}320 & 32.5/45.8/21.6 & OK & 400 ep, fatta \\
    MLP-batch  & 20260623\_171855 & 3180 & 19{.}080 & 35.3/32.8/31.9 & OK & 400 ep, fatta \\
    GNN-curr.\ & 20260630\_181143 & 3293 & 19{.}758 & 32.4/37.4/30.1 & OK & 401 ep, fatta \\
    GNN-batch  & 20260714\_155209 & 3168 & 19{.}008 & 35.4/32.8/31.8 & OK & 400 ep, fatta \\
    \bottomrule
  \end{tabular}
\end{table}

L'esposizione realizzata rispecchia i due design: entrambe le run batch
restano entro $\approx 2$ punti percentuali dall'uniforme $1/3$ per
livello, con tutti e tre i livelli presenti entro i primi venti episodi;
le run curriculum concentrano l'esposizione easy/medium all'inizio e
raggiungono hard solo dopo i loro sblocchi
(\cref{sec:sample-efficiency}), chiudendo con quote hard del 21.6\%
(MLP) e 30.1\% (GNN).

% ------------------------------------------------------------
\section{Efficienza campionaria}
\label{sec:sample-efficiency}

\begin{figure}[t]
  \centering
  \includegraphics[width=\textwidth]{learning_curves}
  \caption{Tasso di successo a finestra (100 episodi) rispetto ai passi
  d'ambiente, tutte e quattro le celle. Avvertenza importante: ogni
  curva è misurata sul \emph{mix di livelli concorrente} di quella run
  --- gated per le celle curriculum, uniforme per batch --- quindi i
  confronti verticali tra bracci non sono omogenei finché il curriculum
  non ha sbloccato tutti i livelli; il confronto a distribuzione
  appaiata è la valutazione finale (\cref{sec:final-performance}).}
  \label{fig:learning-curves}
\end{figure}

\paragraph{Fase iniziale: decollo indistinguibile.}
Tutte e quattro le celle salgono alla stessa velocità: l'SR veicoli a
finestra di 100 episodi supera per la prima volta il 30\% dopo
137--142K passi e il 50\% dopo 162--175K passi in ogni cella --- uno
scarto sotto l'8\%, ben dentro il rumore. La parte notevole è
\emph{su che cosa} ogni policy lo ha superato: le celle curriculum su
esposizione solo-easy, le celle batch sul mix uniforme completo che
include percorsi da \SI{100}{m}. A parità di budget, il decollo del
tasso di successo iniziale non è stato il fattore discriminante tra i
regimi in questa campagna.

\paragraph{Tempi di sblocco (celle curriculum).}
Medium si è sbloccato praticamente nello stesso punto in entrambe le run
curriculum (episodio 842 di 3220, 26.1\%, MLP; episodio 860 di 3293,
26.1\%, GNN). Hard, invece, si è sbloccato molto prima sotto il critic
GNN: episodio 1{.}910 (58.0\% dell'addestramento) contro 2{.}286
(71.0\%) per MLP --- quest'ultimo coerente con lo scatto del cap di
pressione al 70\% del budget più che con il gate di competenza
(\cref{sec:regimes}), mentre la cella GNN ha raggiunto il gate hard per
competenza. Il risultato è visibile nelle quote di esposizione della
\cref{tab:runmatrix}: la policy GNN-curriculum ha visto $\approx 40$\%
episodi hard in più della controparte MLP.

\paragraph{Dopo il decollo.}
Da $\approx 0.5$M passi le curve batch si assestano sotto le due curve
curriculum --- MLP-batch nel modo più visibile
(\cref{fig:learning-curves}) --- che si seguono da vicino; tutte le
curve derivano verso il basso a fine addestramento man mano che i mix
del curriculum si induriscono (celle curriculum) e che il successo a
finestra dei pedoni declina --- le dinamiche entro la run sono
analizzate nella \cref{sec:behavior}.

% ------------------------------------------------------------
\section{Prestazione cumulativa finale}
\label{sec:final-performance}

\begin{table}[t]
  \centering\small
  \caption{Metriche cumulative di addestramento per cella (a livello di
  agente, \%; velocità in km/h per i veicoli, m/s per i pedoni).}
  \label{tab:cumulative}
  \begin{tabular}{@{}lrrrrrrrr@{}}
    \toprule
    & \multicolumn{5}{c}{Veicoli} & \multicolumn{3}{c}{Pedoni} \\
    \cmidrule(lr){2-6}\cmidrule(l){7-9}
    Cella & SR & s$+$t & coll. & off. & vel. & SR & s$+$t & vel. \\
    \midrule
    MLP-curr.\ & 57.28 & 30.48 & 8.84 & 3.41 & 10.2 & 56.59 & 42.56 & 1.96 \\
    MLP-batch  & 46.31 & 23.61 & 17.33 & 12.76 & 18.7 & 57.06 & 42.03 & 1.61 \\
    GNN-curr.\ & 56.55 & 20.84 & 13.37 & 9.23 & 18.0 & 60.50 & 38.58 & 1.91 \\
    GNN-batch  & 49.41 & 22.75 & 17.56 & 10.28 & 26.6 & 58.03 & 41.11 & 1.45 \\
    \bottomrule
  \end{tabular}
\end{table}

\begin{table}[t]
  \centering\small
  \caption{Valutazione finale a 400 episodi, coppia MLP (a livello di
  agente, \%). Scenari: easy/medium/hard su Town03, test su Town05.}
  \label{tab:eval-mlp}
  \begin{tabular}{@{}llrrrrr@{}}
    \toprule
    Scenario & Cella & SR veic. & s$+$t veic. & coll. veic. & off. veic. & SR ped. \\
    \midrule
    easy   & MLP-curr.\ & 60.33 & 28.00 & 7.33 & 4.33 & 58.00 \\
           & MLP-batch  & 41.33 & 28.33 & 18.00 & 12.33 & 58.67 \\
    medium & MLP-curr.\ & 51.67 & 36.00 & 9.00 & 3.33 & 40.67 \\
           & MLP-batch  & 31.00 & 27.67 & 25.33 & 16.00 & 41.67 \\
    hard   & MLP-curr.\ & 44.33 & 38.33 & 14.33 & 3.00 & 26.67 \\
           & MLP-batch  & 29.00 & 29.67 & 23.33 & 18.00 & 32.33 \\
    test   & MLP-curr.\ & 37.00 & 50.67 & 11.00 & 1.33 & 62.67 \\
           & MLP-batch  & 18.00 & 39.00 & 37.33 & 5.67 & 63.00 \\
    \midrule
    tutti  & MLP-curr.\ & 48.33 & 38.25 & 10.42 & 3.00 & 47.00 \\
           & MLP-batch  & 29.83 & 31.17 & 26.00 & 13.00 & 48.92 \\
    \bottomrule
  \end{tabular}
\end{table}

\begin{table}[t]
  \centering\small
  \caption{Valutazione finale, coppia GNN (a livello di agente, \%;
  stesso protocollo della \cref{tab:eval-mlp}; lo scenario easy di
  GNN-curriculum mantiene 101 episodi, \cref{sec:integrity}). Analizzata
  nella \cref{sec:arch-axis}.}
  \label{tab:eval-gnn}
  \begin{tabular}{@{}llrrrrr@{}}
    \toprule
    Scenario & Cella & SR veic. & s$+$t veic. & coll. veic. & off. veic. & SR ped. \\
    \midrule
    easy   & GNN-curr.\ & 52.48 & 22.11 & 13.53 & 11.88 & 71.95 \\
           & GNN-batch  & 55.67 & 22.67 & 14.00 & 7.67 & 63.00 \\
    medium & GNN-curr.\ & 42.33 & 19.67 & 24.67 & 13.33 & 53.67 \\
           & GNN-batch  & 42.00 & 18.67 & 28.33 & 11.00 & 46.33 \\
    hard   & GNN-curr.\ & 36.00 & 25.33 & 25.67 & 13.00 & 39.67 \\
           & GNN-batch  & 33.33 & 22.67 & 28.33 & 15.67 & 31.33 \\
    test   & GNN-curr.\ & 35.67 & 33.33 & 28.00 & 3.00 & 70.67 \\
           & GNN-batch  & 25.67 & 31.33 & 37.00 & 6.00 & 54.00 \\
    \midrule
    tutti  & GNN-curr.\ & 41.65 & 25.10 & 22.94 & 10.31 & 59.02 \\
           & GNN-batch  & 39.17 & 23.83 & 26.92 & 10.08 & 48.67 \\
    \bottomrule
  \end{tabular}
\end{table}

\paragraph{Il contrasto primario, coppia MLP.}
Nella valutazione finale a distribuzione appaiata
(\cref{tab:eval-mlp}), la policy veicolo addestrata a curriculum
completa il 48.33\% dei percorsi contro il 29.83\% del batch: $+18.5$
punti percentuali, con lo stesso segno in \emph{tutti e quattro} gli
scenari ($+19.0$, $+20.7$, $+15.3$, $+19.0$\,pp) e nel cumulativo di
addestramento ($+11.0$\,pp, \cref{tab:cumulative}). Il margine è circa
da due a quattro volte la banda delle repliche a seed identico della
\cref{sec:iterations}, e la sua coerenza tra scenari, livelli di
difficoltà e strati di misura è ciò che il design appaiato può offrire
al posto della replicazione dei seed (\cref{sec:campaign}).

\paragraph{Due modi diversi di fallire.}
I regimi producono guide riconoscibilmente diverse
(\cref{fig:failure-modes}). La policy batch è quasi due volte più
veloce (18.7 contro 10.2 km/h in addestramento) ma converte quella
mobilità in contatto: collisione 17.3\% contro 8.8\% e offroad 12.8\%
contro 3.4\% in addestramento, che si allargano a collisione 26.0\%
contro 10.4\% in valutazione. La policy curriculum fallisce
principalmente per stallo (stuck$+$timeout 38.3\% contro 31.2\% in
valutazione). Nessun profilo domina l'altro su ogni asse --- il
curriculum scambia velocità per sicurezza --- ma sul criterio di
successo rigoroso lo scambio si risolve chiaramente a suo favore.

\begin{figure}[t]
  \centering
  \includegraphics[width=\textwidth]{failure_modes}
  \caption{Composizione delle terminazioni dei veicoli per cella, in
  addestramento (cumulativo) e alla valutazione finale. Le due
  grammatiche del fallimento sono visibili a colpo d'occhio: dominata
  dallo stallo (giallo) in MLP-curriculum, dominata dal contatto
  (rosso/viola) nelle celle batch, con le celle GNN intermedie.}
  \label{fig:failure-modes}
\end{figure}

\paragraph{Pedoni: nessun effetto di regime.}
Il successo dei pedoni è piatto tra i bracci ovunque: 56.6 contro
57.1\% in addestramento, 47.0 contro 48.9\% in valutazione, ed entro
$\pm 6$\,pp in ogni scenario. Il regime di addestramento, in questa
piattaforma, è una variabile della policy veicolo; il task pedone appare
insensibile all'ordinamento dei task. La struttura dei fallimenti dei
pedoni è ripresa nella \cref{sec:ped-limits}.

\paragraph{Le celle GNN a colpo d'occhio.}
In addestramento, la cella GNN-curriculum atterra entro il rumore di
MLP-curriculum sull'SR principale (56.55 contro 57.28\%)
ridistribuendo però la massa dei fallimenti: meno stallo ($-9.6$\,pp
stuck$+$timeout), più contatto ($+4.5$\,pp collisione, $+5.8$\,pp
offroad) e velocità molto più alta (18.0 contro 10.2 km/h) --- più
vicina alla mobilità delle celle batch con successo da curriculum.
GNN-batch spinge le stesse tendenze oltre (26.6 km/h, la cella più
veloce) \emph{battendo} al contempo MLP-batch sull'SR di addestramento
di $+3.1$\,pp e in valutazione di $+9.3$\,pp (\cref{tab:eval-gnn}).
L'encoder del critic, dunque, plasma il profilo di rischio e
interagisce con il regime; la lettura architetturale completa è nella
\cref{sec:arch-axis}.

% Fonti: campaign_consolidation.json + analisi ad-hoc per livello×quartile
% (ricalcolata 2026-07-16; le traiettorie Q easy controllano il mix).
\section{Analisi comportamentale}
\label{sec:behavior}

\subsection{Dinamica entro la run: erosione, non solo indurimento del mix}
\label{sec:erosion}

Ogni curva di apprendimento della \cref{fig:learning-curves} deriva
verso il basso dopo il suo picco iniziale. Per le celle curriculum
questo è in parte meccanico --- il loro mix concorrente si indurisce a
ogni sblocco --- quindi la \cref{tab:quartiles} separa i due effetti:
il pannello superiore riporta le traiettorie per quartile complessive,
quello inferiore la traiettoria \emph{sui soli episodi easy}, dove il
task è costante per costruzione.

\begin{table}[t]
  \centering\small
  \caption{SR veicoli per quartile di addestramento (\%). Sopra: tutti
  i livelli (confuso dal mix per le celle curriculum). Sotto: solo
  episodi easy (controllato per livello; easy resta $\approx 30$\% del
  mix dopo gli sblocchi). $\Delta$ = Q4 $-$ Q1.}
  \label{tab:quartiles}
  \begin{tabular}{@{}lrrrrr@{}}
    \toprule
    Cella & Q1 & Q2 & Q3 & Q4 & $\Delta$ \\
    \midrule
    \multicolumn{6}{@{}l}{\emph{Tutti i livelli}} \\
    MLP-curr.\ & 68.9 & 61.1 & 54.4 & 44.7 & $-24.2$ \\
    MLP-batch  & 58.7 & 45.9 & 41.8 & 38.7 & $-20.0$ \\
    GNN-curr.\ & 66.6 & 62.2 & 54.5 & 42.9 & $-23.7$ \\
    GNN-batch  & 50.5 & 54.6 & 47.5 & 45.1 & $-5.4$ \\
    \midrule
    \multicolumn{6}{@{}l}{\emph{Solo livello easy}} \\
    MLP-curr.\ & 68.9 & 69.3 & 59.6 & 52.4 & $-16.5$ \\
    MLP-batch  & 64.2 & 52.0 & 48.9 & 45.6 & $-18.6$ \\
    GNN-curr.\ & 66.6 & 67.9 & 56.1 & 55.8 & $-10.8$ \\
    GNN-batch  & 59.7 & 62.8 & 54.3 & 53.3 & $-6.4$ \\
    \bottomrule
  \end{tabular}
\end{table}

Il pannello controllato per livello mostra che l'indurimento del mix
\emph{non} è tutta la storia: la prestazione a livello easy si erode
entro ogni run, da $-6.4$\,pp (GNN-batch, la cella più stabile) a
$-18.6$\,pp (MLP-batch). Le celle batch rendono il caso in modo pulito
--- la loro distribuzione dei task è costante, quindi il loro declino è
interamente degradazione entro la run --- e le celle curriculum lo
riproducono sulla loro fetta easy. I pedoni si erodono ancora di più
(SR pedoni a livello easy: $-31.8$\,pp in MLP-curriculum, da $-16$ a
$-22$\,pp altrove). Due meccanismi candidati sono compatibili con i
dati e non separabili a $n{=}1$: la co-adattazione tra classi (il
comportamento dei pedoni degrada in parallelo e alimenta il canale di
hazard dei veicoli, \cref{sec:obs-actions}) e lo schedule di annealing
dell'entropia --- benché la transizione si completi solo all'83\% del
budget, dentro Q4, mentre l'erosione inizia già in Q3, quindi
l'annealing non può esserne l'unica causa. La conclusione pulita è
architetturale: sotto il critic GNN l'erosione è dimezzata o meglio,
qualunque ne sia il driver (\cref{sec:arch-axis}).

\subsection{Quattro policy, due grammatiche del fallimento, un'eccezione}
\label{sec:failure-grammars}

La velocità media ordina le quattro celle: 10.2 (MLP-curr.), 18.0
(GNN-curr.), 18.7 (MLP-batch), 26.6 km/h (GNN-batch); i fallimenti da
contatto seguono lo stesso ordine (collisione in addestramento 8.8 /
13.4 / 17.3 / 17.6\%), mentre lo stallo anti-correla (stuck$+$timeout
30.5 / 20.8 / 23.6 / 22.8\%). Ciò che \emph{non} varia è la struttura
interna dello stallo: in tutte e quattro le celle, le terminazioni
stuck si decompongono quasi identicamente ($44$--$47$\% rilevatore di
loop più non-progresso a basso completamento, $39$--$42$\% loop più
non-progresso a completamento più alto, il resto minore), e i passi
medi di non-progresso coprono una banda stretta 216--253. Regimi e
critic cambiano \emph{quanto} una policy stalla, non \emph{come} stalla
--- coerente con l'idea che lo stallo sia un attrattore unico
(l'equilibrio difensivo della \cref{sec:iterations}) da cui
l'addestramento o esce o non esce.

\subsection{L'equilibrio difensivo non è intrinseco al curriculum}
\label{sec:defensive-case}

Il profilo lento-e-in-stallo appartiene a esattamente una cella:
MLP-curriculum (10.2 km/h, stuck$+$timeout 30.5\%). La sua gemella
curriculum sotto il critic GNN è quasi due volte più veloce (18.0 km/h)
con un terzo di stallo in meno (20.8\%) allo stesso SR principale ---
quindi l'equilibrio cauto documentato durante lo sviluppo (il gate di
hazard che tiene fermi i veicoli, \cref{sec:iterations}) non è una
proprietà dell'addestramento a curriculum in quanto tale. È un bacino
in cui la run con critic MLP si è assestata e che la run con critic GNN
ha evitato, plausibilmente perché la stima del valore su input
strutturati per agente prezza diversamente il canale di hazard. A
$n{=}1$ per cella questa resta un'interpretazione; ciò che i dati
stabiliscono è che il regime da solo non determina l'equilibrio.

\section{Generalizzazione: da Town03 a Town05}
\label{sec:generalization}

Lo scenario test colloca le stesse policy su Town05 --- un layout a
griglia quadrata e più corsie, strutturalmente diverso da Town03
(\cref{sec:carla-bg}) --- per 100 episodi con target di \SI{80}{m} e
traffico come in addestramento. Tre osservazioni dalla coppia MLP
(\cref{tab:eval-mlp}) e una quarta dalla coppia GNN
(\cref{tab:eval-gnn}); la \cref{fig:town05} riassume il quadro del
trasferimento.

\begin{figure}[t]
  \centering
  \includegraphics[width=0.8\textwidth]{town05_bars}
  \caption{SR veicoli: media degli scenari Town03 vs.\ test Town05, per
  cella. Le celle curriculum calano meno; la coppia GNN è
  indistinguibile in-domain (43.6 vs.\ 43.7\%) e si separa di 10\,pp
  sulla mappa mai vista.}
  \label{fig:town05}
\end{figure}

\paragraph{L'ordinamento sopravvive al cambio di mappa.}
Il curriculum completa il 37.0\% dei percorsi su Town05 contro il
18.0\% del batch ($+19.0$\,pp) --- essenzialmente lo stesso margine
in-domain. Il calo assoluto dalla media degli scenari Town03 è quasi
identico per i due bracci ($-15.1$\,pp da 52.1\% per il curriculum;
$-15.8$\,pp da 33.8\% per il batch), quindi il vantaggio del curriculum
non è un'idiosincrasia di Town03: si trasferisce per intero su strade
mai viste in addestramento. Letto contro l'ipotesi H2, la copertura
iniziale più ampia del regime batch non gli ha comprato \emph{alcun}
vantaggio di generalizzazione --- semmai il calo relativo è peggiore
dalla sua base più bassa ($-47$\% contro $-29$\% relativo).

\paragraph{Le firme di fallimento si amplificano fuori distribuzione.}
Ogni braccio fallisce su Town05 nel modo in cui fallisce a casa, solo
di più: la policy curriculum stalla (stuck$+$timeout 50.7\%, il suo
massimo ovunque) mantenendo basso il contatto (collisione 11.0\%,
offroad 1.3\%); la policy batch si schianta (collisione 37.3\%, il suo
massimo ovunque). Le strade nuove stressano esattamente l'asse su cui
ogni policy già poggiava.

\paragraph{Sotto il critic GNN, il vantaggio esiste \emph{solo} fuori
distribuzione.}
La coppia GNN affila il punto. In-domain i due bracci GNN sono gemelli
statistici --- SR veicoli medio sugli scenari Town03 43.6\%
(curriculum) contro 43.7\% (batch) --- eppure su Town05 si separano di
$+10.0$\,pp (35.7 contro 25.7\%), con il calo su Town05 più piccolo di
ogni cella per GNN-curriculum ($-7.9$\,pp) contro $-18.0$\,pp per
GNN-batch. Dove la coppia MLP mostrava un vantaggio del curriculum che
si trasferisce, la coppia GNN ne mostra uno che \emph{compare} solo nel
trasferimento: evidenza con lo stesso segno che l'ordinamento dei task
compra robustezza alle strade nuove, veicolata da due firme diverse.
L'interazione è analizzata nella \cref{sec:arch-axis}.

\paragraph{I pedoni si trasferiscono con facilità.}
Il successo dei pedoni su Town05 (62.7 contro 63.0\%) supera i
risultati hard di entrambi i bracci su Town03 di $\approx 30$\,pp ---
i percorsi test da \SI{80}{m} sono più fattibili dei percorsi hard da
\SI{100}{m} di Town03, indicando ancora la fattibilità dei percorsi,
non la generalizzazione della policy, come collo di bottiglia dei
pedoni (\cref{sec:ped-limits}).

% Fonti: campaign_consolidation.json; delta ricalcolati 2026-07-16.
\section{Asse architetturale: critic MLP vs GNN}
\label{sec:arch-axis}

L'asse secondario sostituisce solo l'encoder del critic
(\cref{sec:critic-encoders}); actor, reward, ambiente e budget sono
identici, quindi le differenze tra le colonne MLP e GNN delle tabelle
precedenti misurano come la \emph{stima del valore} plasma ciò che gli
stessi actor apprendono. Ne seguono due letture: l'effetto proprio
dell'architettura entro ogni regime e --- la domanda H3 --- se
l'effetto di regime sopravviva al cambio di critic.

\paragraph{L'effetto di regime dipende dal critic.}
La \cref{tab:interaction} condensa il risultato centrale della
campagna. Sotto il critic MLP, il curriculum batte il batch con un
margine ampio e coerente tra gli strati. Sotto il critic GNN lo stesso
contrasto quasi svanisce: $+2.5$\,pp complessivi in valutazione ---
dentro la banda di rumore a seed identico della \cref{sec:iterations}
--- ed esattamente zero in-domain, dove i due bracci GNN sono gemelli
statistici (media scenari Town03 43.60 contro 43.67\%). L'unica
componente del vantaggio del curriculum che \emph{si replica in
entrambe le architetture} è quella out-of-distribution: $+19.0$\,pp
(MLP) e $+10.0$\,pp (GNN) su Town05. L'ipotesi H3 è quindi supportata
nella sua forma forte: l'encoder del critic cambia non solo l'entità ma
il luogo stesso dell'effetto di regime.

\begin{table}[t]
  \centering\small
  \caption{L'interazione arch$\times$regime: SR veicoli
  curriculum-meno-batch (pp) per strato di misura, entro ciascuna
  architettura. Delta di valutazione dei pedoni aggiunto per
  completezza.}
  \label{tab:interaction}
  \begin{tabular}{@{}lrr@{}}
    \toprule
    $\Delta$(curr.\ $-$ batch) & Coppia MLP & Coppia GNN \\
    \midrule
    Cumulativo di addestramento (veic.) & $+10.97$ & $+7.14$ \\
    Valutazione complessiva (veic.) & $+18.50$ & $+2.48$ \\
    Valutazione, media Town03 (veic.) & $+18.33$ & $-0.06$ \\
    Valutazione, Town05 (veic.) & $+19.00$ & $+10.00$ \\
    Valutazione complessiva (ped.) & $-1.92$ & $+10.35$ \\
    \bottomrule
  \end{tabular}
\end{table}

\paragraph{Da dove viene l'interazione.}
La compressione è prodotta da due effetti architetturali asimmetrici,
uno per regime. Entro il batch, il critic GNN \emph{aiuta}: $+3.1$\,pp
in addestramento e $+9.3$\,pp in valutazione rispetto a MLP-batch,
insieme alla dinamica entro la run più stabile della campagna (l'unica
cella senza declino monotono dei quartili e con la più piccola erosione
controllata per livello, \cref{tab:quartiles}). È coerente con il
meccanismo che motiva i critic CTDE in primo luogo (\cref{sec:marl}):
sotto una distribuzione mista con sei agenti interagenti, la stima
strutturata del valore per agente dovrebbe aiutare di più dove il
segnale di apprendimento è più rumoroso --- cioè il regime batch dal
passo~0. Entro il curriculum, il critic GNN rimodella più che
migliorare: evita l'equilibrio difensivo (\cref{sec:defensive-case}),
sblocca hard per competenza anziché per cap di pressione
(\cref{sec:sample-efficiency}), vede $\approx 40$\% episodi hard in più
--- e converte tutto questo in mobilità e contatto anziché in successo
(collisione in valutazione 22.9 contro 10.4\%, SR $-6.7$\,pp rispetto a
MLP-curriculum). I pedoni mostrano l'interazione speculare: il critic
GNN li solleva sotto curriculum ($+12.0$\,pp in valutazione, la miglior
cella pedoni della campagna al 59.0\%) ma non sotto batch
($-0.3$\,pp).

\paragraph{La firma di aprile non ricompare.}
La lettura qualitativa del pilota (\cref{sec:pilot}) ---
GNN-curriculum come la cella più \emph{difensiva}, stallo massimo con
contatto minimo --- si inverte sulla piattaforma sana: l'attuale
GNN-curriculum è la cella curriculum \emph{meno} in stallo, con i tassi
di contatto più alti. Una firma qualitativa misurata a SR di pavimento,
in un ambiente pre-correzioni, non è sopravvissuta alle correzioni
dell'ambiente --- un'altra istanza della regola di comparabilità tra
ere della \cref{sec:bugs}, e un monito contro il leggere effetti di
architettura da distribuzioni di task rotte.

\paragraph{Avvertenze.}
Tutto quanto sopra poggia su un contrasto appaiato per coppia di celle:
il pattern di interazione è internamente coerente su quattro strati di
misura ed entrambe le classi di agenti, ma le sue grandezze portano la
consueta qualifica $n{=}1$ (\cref{sec:threats}). Il delta complessivo
di $+2.5$\,pp della coppia GNN esplicitamente \emph{non} è rivendicato
come vantaggio del curriculum; il risultato rivendicato è il suo
contrasto con la coppia MLP e la componente Town05 replicata tra le
architetture.

% Fonti: campaign_consolidation.json (diagnostica percorsi ped. + SR
% per livello, tutte e quattro le celle, ricalcolato 2026-07-16).
\section{Limiti strutturali dei pedoni}
\label{sec:ped-limits}

I risultati dei pedoni vanno letti contro un vincolo che nessun regime
di addestramento può rimuovere: la topologia dei marciapiedi di Town03
spesso non può fornire il task nominale.

\paragraph{Le diagnostiche sono invarianti per regime e architettura.}
In tutte e quattro le celle le diagnostiche di generazione dei percorsi
pedonali sono statisticamente indistinguibili --- percorsi under-target
55.1--55.9\% dei record, flag too-short 7.9--9.2\%, percorsi da respawn
13.4--14.1\%, e sorgenti divise in $\approx$40\% catene con
attraversamenti / $\approx$46\% catene di solo marciapiede. Questa
invarianza è il punto: questi numeri sono proprietà della mappa e del
planner (\cref{sec:routes}), non delle policy, e fissano quindi un
soffitto comune per ogni cella. La guardia di retrocessione
\texttt{route\_short}, al contrario, scatta quasi mai (0.6\% dei record
pedoni): i percorsi degeneri corti sono rari; i percorsi \emph{più
corti del nominale ma sopra la guardia} sono la norma.

\paragraph{Il successo collassa con la distanza target --- ovunque.}
L'SR pedoni in addestramento cala monotonicamente con il livello in
ogni cella (easy 70--83\%, medium 53--60\%, hard 29--44\%), e la
valutazione a distribuzione appaiata riproduce il gradiente (easy
58--72\%, medium 41--54\%, hard 27--40\%;
\cref{tab:eval-mlp,tab:eval-gnn}). Un collasso che sopravvive a due
architetture, due regimi e due strati di misura, mentre le diagnostiche
dei percorsi restano costanti, è guidato dalla geometria del task: a
target nominali di \SI{100}{m}, più della metà dei percorsi generati
non può nemmeno raggiungere la lunghezza target, e quelli che la
raggiungono richiedono lunghe catene di marciapiede fattibili che
Town03 spesso non ha.

\paragraph{Il controfattuale Town05.}
Lo scenario test rende concreta l'attribuzione: su Town05 con target di
\SI{80}{m}, l'SR pedoni balza a 54--71\% in ogni cella --- da 15 a
31\,pp \emph{sopra} il risultato hard su Town03 di ciascuna cella, su
una mappa mai vista in addestramento. Se il collasso medium/hard fosse
un limite di policy, strade mai viste dovrebbero nuocere, non aiutare;
invece una topologia di marciapiedi più trattabile recupera gran parte
della perdita. Il vincolo stringente è la fattibilità, non la
generalizzazione.

\paragraph{Che cosa resta attribuibile alla policy.}
Entro l'inviluppo di fattibilità, le differenze di policy sono visibili
e coerenti con la \cref{sec:final-performance}: il passo di camminata
copre 1.45--1.96 m/s tra le celle (tutte dentro la banda di comfort),
le collisioni dei pedoni sono trascurabili ovunque (0.2--0.3\% --- i
pedoni falliscono per tempo e percorso, non per contatto), e una cella
si separa dal gruppo in valutazione (GNN-curriculum, SR pedoni 59.0\%
contro 47--49\% altrove), un filo ripreso nella \cref{sec:arch-axis}.
La conseguenza metodologica resta comunque: i tassi di successo
medium/hard dei pedoni non sono una misura pulita della qualità della
policy su questa piattaforma --- questa tesi confronta quindi i pedoni
solo entro livello e tra bracci, e la riparazione della generazione dei
percorsi è elencata come lavoro futuro (\cref{ch:conclusions}).

\section{Risultati rispetto alle ipotesi}
\label{sec:findings}

La \cref{tab:hypotheses} chiude il capitolo rispetto alle ipotesi della
\cref{sec:rq}. Due cose sopravvivono a ogni cella, ogni strato ed
entrambe le architetture: il vantaggio del curriculum sulla mappa mai
vista, e il ruolo del regime come variabile della policy \emph{veicolo}
che rimodella la composizione dei fallimenti. Ciò che non è
sopravvissuto è l'aspettativa popolare da entrambi i lati: il
curriculum non ha appreso più in fretta all'inizio, e la copertura più
ampia del batch non ha comprato robustezza.

\begin{table}[t]
  \centering\small
  \caption{Ipotesi contro evidenza (contrasti appaiati, $n{=}1$ per
  cella; delta in pp di SR veicoli salvo diversa indicazione).}
  \label{tab:hypotheses}
  \begin{tabular}{@{}p{2.2cm}p{4.6cm}p{2.2cm}p{4.0cm}@{}}
    \toprule
    Ipotesi & Predizione & Verdetto & Evidenza \\
    \midrule
    H1 (curriculum) & migliore efficienza campionaria \emph{iniziale} & rigettata &
      decollo identico in tutte e quattro le celle (50\% a finestra a
      162--175K passi; \cref{sec:sample-efficiency}) \\
    & prestazione finale più alta & supportata sotto MLP; non sotto GNN
      in-domain &
      $+18.5$ (MLP) vs.\ $+2.5$ (GNN) in valutazione;
      \cref{tab:interaction} \\
    & rischio di plateau / oblio su hard & rimodellato &
      l'erosione tardiva è reale ma \emph{di piattaforma}, anche nelle
      celle batch; $\Delta$Q1$\to$Q4 controllato per livello da $-6$ a
      $-19$\,pp (\cref{tab:quartiles}) \\
    \addlinespace
    H2 (batch) & piena copertura dal passo 0 & confermata (meccanica) &
      tutti i livelli entro i primi 20 episodi; mix realizzato
      $\approx$ uniforme (\cref{tab:runmatrix}) \\
    & segnale iniziale più rumoroso & non visibile &
      stesso decollo del curriculum sul mix più difficile
      (\cref{sec:sample-efficiency}) \\
    & copertura $\to$ robustezza su mappa nuova & rigettata &
      calo su Town05 uguale o peggiore del curriculum in entrambe le
      coppie ($-15.8$/$-18.0$ vs.\ $-15.1$/$-7.9$\,pp;
      \cref{sec:generalization}) \\
    \addlinespace
    H3 (interazione) & l'effetto di regime può dipendere dall'encoder
      del critic &
      supportata (forma forte) &
      effetto in-domain $+18.3$ (MLP) vs.\ $-0.1$ (GNN); solo la
      componente Town05 si replica ($+19.0$/$+10.0$;
      \cref{tab:interaction}) \\
    \bottomrule
  \end{tabular}
\end{table}

Formulato come risposta alla domanda di ricerca: l'ordinamento dei task
\emph{produce} un comportamento di guida multi-agente misurabilmente
diverso a parità di budget --- ma non dove i curricula vengono di
solito pubblicizzati. Non ha accelerato l'apprendimento iniziale; ha
cambiato ciò a cui la policy veicolo è convergiuta (cauta-e-in-stallo
contro veloce-e-in-collisione), ha lasciato i pedoni essenzialmente
intatti, il suo vantaggio di successo in-domain esiste solo sotto il
critic MLP, e il suo beneficio più portabile --- coerente in entrambe
le architetture del critic --- è la robustezza su strade mai viste in
addestramento. Ogni affermazione quantitativa qui sopra eredita la
qualifica $n{=}1$-per-cella della \cref{sec:campaign};
l'interpretazione di questi risultati è ripresa nel
\cref{ch:discussion}.
